\begin{proof}[Proof of Lemma~\ref{sga3:XXI:5.2}, continued]
Put \(m=n_{ij}\), \(s_i=s_{\alpha_i}\), and
\(s_j=s_{\alpha_j}\). We have
\[
 T_jT_i=(T_iT_j)^{m-1}.
\]
By Corollary~\ref{sga3:XXI:4.4},
\(\operatorname{ord}_\Delta(w(\alpha))<p\) for every
\(w\in W_{ij}\) distinct from
\[
 1=(s_is_j)^m
 \quad\text{and}\quad
 s_j=s_i(s_is_j)^{m-1}.
\]
The induction hypothesis therefore gives
\[
\begin{split}
 &T_j(T_iT_j)^{m-2}T_{s_i(\alpha)}
       (T_jT_i)^{m-2}T_j\\
 &\qquad
 =T_{s_j(s_is_j)^{m-2}s_i(\alpha)}
 =T_{s_is_j(\alpha)}
 =T_{s_i(\alpha)}.
\end{split}
\]
The last equality holds because \(s_j(\alpha)=\alpha\). Hence
\[
 (T_iT_j)^{m-1}T_{s_i(\alpha)}(T_jT_i)^{m-1}
 =T_iT_{s_i(\alpha)}T_i,
\]
which proves \eqref{sga3:XXI:5.2:orthogonal-expanded}.
\end{proof}

\begin{lemma}\label{sga3:XXI:5.3}
Let \(h\in\widetilde W\). Write
\[
 h=T_{\alpha_1}\cdots T_{\alpha_m},
\]
where the \(\alpha_i\in\Delta\), not necessarily distinct, are chosen
so that \(m\) is minimal. Then
\[
 p(h)(\alpha_m)\in-\mathscr{P}(\Delta).
\]
\end{lemma}

\begin{proof}
Since
\[
 p(T_{\alpha_m})(\alpha_m)
 =s_{\alpha_m}(\alpha_m)=-\alpha_m,
\]
if \(p(h)(\alpha_m)\) were positive, there would be an index
\(k\), \(1\le k\le m-1\), such that
\[
 u=s_{\alpha_{k+1}}\cdots s_{\alpha_m}(\alpha_m)
 =-s_{\alpha_{k+1}}\cdots s_{\alpha_{m-1}}(\alpha_m)
 \in-\mathscr{P}(\Delta)
\]
and \(s_{\alpha_k}(u)\in\mathscr{P}(\Delta)\). Necessarily
\(u=-\alpha_k\), by Lemma~\ref{sga3:XXI:3.3.1}, whence
\[
 s_{\alpha_{k+1}}\cdots s_{\alpha_{m-1}}(\alpha_m)=\alpha_k.
\]
Property \SGARefLink{sga3:XXI:3.6:item:iii}{(iii)} of Lemma~\ref{sga3:XXI:5.2} then gives
\[
 T_{\alpha_k}T_{\alpha_{k+1}}\cdots
 T_{\alpha_{m-1}}T_{\alpha_m}
 =
 T_{\alpha_{k+1}}\cdots T_{\alpha_{m-1}},
\]
contradicting the minimality of \(m\).
\end{proof}

Let now \(h\in\widetilde W\) satisfy
\[
 p(h)\bigl(\mathscr{P}(\Delta)\bigr)
 \subset\mathscr{P}(\Delta).
\]
Lemma~\ref{sga3:XXI:5.3} gives \(p(h)=1\). This proves
Theorem~\ref{sga3:XXI:5.1} and also the following two corollaries.

\begin{corollary}\label{sga3:XXI:5.4}
If \(R_+\) is a system of positive roots and \(w\in W\) satisfies
\(w(R_+)=R_+\), then \(w=1\).
\end{corollary}

\begin{corollary}\label{sga3:XXI:5.5}
The Weyl group acts freely and transitively on the set of systems of
positive roots, respectively on the set of systems of simple roots,
respectively on the set of Weyl chambers.
\end{corollary}

Choose a system of simple roots \(\Delta\), and put
\[
 R^+=\mathscr{P}(\Delta).
\]
\footnote{Editors' note: The superscript notation \(R^+\), rather
than \(R_+\), is used here so that
\(R^+_{\alpha,\beta}=R_{\alpha,\beta}\cap R^+\).}
For every pair of simple roots
\((\alpha,\beta)\in\Delta\times\Delta\), let
\(R_{\alpha,\beta}\) be the set of roots that are linear combinations
of \(\alpha\) and \(\beta\). Put
\[
 R^+_{\alpha,\beta}=R^+\cap R_{\alpha,\beta},
\]
and let \(W_{\alpha,\beta}\) be the Weyl group of
\(R_{\alpha,\beta}\), that is, the subgroup of \(W\) generated by
\(s_\alpha\) and \(s_\beta\).

\begin{theorem}[Tits]\label{sga3:XXI:5.6}
Let \(\alpha\) and \(\beta\) be two simple roots, and let \(w\in W\)
satisfy \(w(\alpha)=\beta\). There exist a sequence of simple roots
\(\alpha_0,\ldots,\alpha_m\) and a sequence
\(w_0,\ldots,w_{m-1}\) of elements of \(W\) satisfying:
\begin{enumerate}
\item[(i)] \(\alpha_0=\alpha\) and \(\alpha_m=\beta\);
\item[(ii)] \(w=w_{m-1}w_{m-2}\cdots w_0\);
\item[(iii)] \(w_i(\alpha_i)=\alpha_{i+1}\) for
  \(0\le i\le m-1\);
\item[(iv)] for every \(i\), \(0\le i\le m-1\), such that
  \(\alpha_i\ne\alpha_{i+1}\), one has
  \(w_i\in W_{\alpha_i,\alpha_{i+1}}\);
\item[(v)] for every \(i\), \(0\le i\le m-1\), such that
  \(\alpha_i=\alpha_{i+1}\), there exists a simple root \(\beta_i\)
  such that \(w_i\in W_{\alpha_i,\beta_i}\).
\end{enumerate}
\end{theorem}
